\subsection{Calculus with Vector Functions}

\content{
\begin{definition} If $\vec{r}(t)$ is a vector function, it's derivative is
defined by 
$$ \vec{r}'(t) = \lim_{h\rightarrow 0} \frac{\vec{r}(t+h)-\vec{r}(t)}{h}. $$
Provided that this limit exists. 
\end{definition}
}{% presentation
}{% nopresentation
}{% comments
}

\content{
\begin{theorem} If $\vec{r}(t) = \la f(t), g(t), h(t)\ra$, and $f$, $g$, and $h$
are differentiable, then 
$$\vec{r}'(t) = \la f'(t), g'(t), h'(t)\ra. $$
\end{theorem}

\begin{definition} The unit vector with the same direction as the tangent vector
to $\vec{r}(t)$ is called the unit tangent vector and is given by 
$$ \vec{T}(t) = \frac{\vec{r}(t)}{|\vec{r}(t)|}. $$
\end{definition}
}{% presentation
}{% nopresentation
}{% comments
}

\content{
\begin{example} Find the derivative of $\vec{r}(t) = \la 1+t^3, te^{-t}, \sin
2t\ra,$ and find the unit tangent vector at the point $t=0$.
\end{example}
}{% presentation
\vspace{2in}
}{% nopresentation
}{% comments
}

\content{
\begin{theorem} Suppose $\vec{u}$ and $\vec{v}$ are differentiable vector valued
functions, and $c$ is a scalar, and $f(t)$ is a real valued function, then 
\begin{itemize}
\item $\frac{d}{dt}(\vec{u}+\vec{v}) = \vec{u}'(t) + \vec{v}'(t)$
\item $\frac{d}{dt}(c\vec{u}(t)) = c\vec{u}'(t)$
\item $\frac{d}{dt}(f(t)\vec{u}(t)) = f(t)\vec{u}'(t) + f'(t)\vec{u}(t)$
\item $\frac{d}{dt}(\vec{u}(t)\cdot \vec{v}(t)) = \vec{u}'(t)\cdot \vec{v}(t) +
\vec{u}(t)\cdot \vec{v}'(t)$
\item $\frac{d}{dt}(\vec{u}(t)\times\vec{v}(t)) = \vec{u}'(t)\times\vec{v}(t) +
\vec{u}(t)\times\vec{v}'(t)$
\item $\frac{d}{dt}(\vec{u}(f(t))) = f'(t)\vec{u}'(f(t))$
\end{itemize}
\end{theorem}
}{% presentation
}{% nopresentation
}{% comments
}

\content{
\begin{theorem} If $|\vec{r}(t)|=c$ is a constant for all $t$, then
$\vec{r}'(t)$ is orthogonal to $\vec{r}(t)$ for all $t$. 
\end{theorem}
}{% presentation
\vspace{2in}
}{% nopresentation
}{% comments
Start with $\vec{r}(t)\cdot\vec{r}(t) = c^2$, then use the previous theorem to
compute $\frac{d}{dt}(\vec{r}(t)\cdot \vec{r}(t))$.
}

\content{
\begin{example} Compute the derivative of the vector valued function defined by 
$$ \vec{r}(t) = \sin(t)\cdot \la 1,t, e^t\ra. $$
\end{example}
}{% presentation
\vspace{2in}
}{% nopresentation
}{% comments
}

\content{
The integral of a vector valued function is similarly computed: If
$\vec{r}(t)=\la f(t), g(t), h(t)\ra$, and each of $f$, $g$, and $h$ is
integrable on $[a,b]$, then 
$$ \int_a^b \vec{r}(t) dt = \la \int_a^b f(t)dt, \int_a^b g(t)dt, \int_a^b
h(t)dt \ra. $$
}{% presentation
}{% nopresentation
}{% comments
}

\content{
\begin{example}
Compute the indefinite integral of the vector valued function 
$$ \vec{r}(t) = \la 2\cos t, \sin t, 2t\ra. $$
\end{example}
}{% presentation
\vspace{2in}
}{% nopresentation
}{% comments
}
\subsection{Arc Length}


\content{
\begin{definition} Given a vector valued fuction $\vec{r}(t)$, defined on an
interval $[a,b]$, the length of the curve $\vec{r}(t)$ from $a$ to $b$ is given
by 
$$ L = \int_a^b |\vec{r}'(t)|\, dt. $$
\end{definition}
}{% presentation
}{% nopresentation
}{% comments
}

\content{
\begin{example} Find the length of the curve $\vec{r}(t) = \la \cos t, \sin t,
t\ra$ from $a=0$ to $b=2\pi$.
\end{example}
}{% presentation
\vspace{2in}
}{% nopresentation
}{% comments
}

\subsubsection{The Arc Length Function}

\content{
\begin{definition}
Let $\vec{r}(t)$ be a vector valued function, defined on $[a,b]$, we can define
it's arc length function as 
$$ s(t) = \int_a^t |\vec{r}(u)|\, du. $$
\end{definition}
}{% presentation
\vspace{3in}
}{% nopresentation
}{% comments
Note that by the fundamental Theorem of Calculus, we have that 
$\frac{ds}{dt} = |\vec{r}(t)|$. 

It will sometimes be usefull to reparametrize the function $\vec{r}(t)$ in terms
of it's arc length: if $s(t)$ is the arc length function, we might be able to
solve for $t$, that is find a function $t=t(s)$, then the point $\vec{r}(t(3))$
is the point where the particle has traveled 3 units of length from it's
starting point. 
}

\subsubsection{The Tangent, Normal, and Binormal Vectors}

\content{
\begin{definition} Given a vector valued function $\vec{r}(t)$ the unit tangent
vector is given by 
$$ \vec{T}(t) = \frac{\vec{r}'(t)}{|\vec{r}'(t)|}. $$
\end{definition}
}{% presentation
\vspace{1in}
}{% nopresentation
}{% comments
This is simply a unit vector in the same direction as the tangent vector to
$\vec{r}(t)$.
}

\content{
\begin{definition}
Given a vector valued function $\vec{r}(t)$, the unit normal vector is given by 
$$ \vec{N}(t) = \frac{\vec{T}'(t)}{|\vec{T}'(t)|}. $$
\end{definition}
}{% presentation
\vspace{2in}
}{% nopresentation
}{% comments
Since $|\vec{T}(t)|=1$ for all $t$, $\vec{T}'(t)$ is orthogonal to $\vec{T}(t)$.
Thus this vector can be thought of as the direction that the particle is turning
at time $t$.
}

\content{
\begin{definition}
Given a vector valued function $\vec{r}(t)$, the binormal vector $\vec{B}(t)$ is
given by 
$$ \vec{B}(t) = \vec{T}(t)\times \vec{N}(t). $$
\end{definition}
}{% presentation
\vspace{2in}
}{% nopresentation
}{% comments
This vector is orthogonal to both $\vec{T}(t)$ and $\vec{N}(t)$, and is also a
unit vector. 

Draw a picture of how this gives a basis (called the TNB basis), which is
important to differential geometry and applications to the motion of spacecraft. 
}

\content{
\begin{example} Find the unit normal vector and binormal vectors for the
function 
$$ \vec{r}(t) = \la \cos t, \sin t, t\ra. $$
\end{example}
}{% presentation
\vspace{5in}
}{% nopresentation
}{% comments
}
